% Supplement to routing_theory_notes.tex; a LaTeX fragment for incorporation.
% Independent derivation reviewed by exhaustive through-partition tests.

\section{Polynomial reblocking for every tile with through components}
\label{sec:through-reblocking}

The alignment hypothesis can be obtained after a bounded amount of explicit
reblocking.  The result below applies whenever every connected component of
the tile meets both interfaces.  It does not require the input and output
partitions to agree initially, and it does not require equal input and output
widths within an individual component.

At the level of boundary connectivity, these tiles are the classical block
bijections, or elements of the dual symmetric inverse monoid
\cite{fitzgerald-leech1998,easdown-east-fitzgerald2008}.  Their composition by
stacking diagrams is exactly the connectivity operation used here.  The
following elementary stabilization proof is included in full; the contribution
for this application is its combination with marked-cut repetition kernels
and source-address provenance.  No priority claim is made for the underlying
semigroup stabilization principle.

\subsection{Stabilization of the component partitions}

Call a tile \emph{through} if every connected component meets both its input
and its output interface.  For such a tile $T$, let $d$ be its number of
components, and write $A_j,B_j$ for their input and output index sets as before.
All $A_j$ and $B_j$ are nonempty, but the unordered partitions need not agree.
Every power $T^{\circ n}$ is again through: each component occurrence can be
continued through a nonempty interface in both directions until the exposed
interfaces of the whole power are reached.

Define a directed seam relation on the original component indices by
\[
 j\mathrel D l \quad\Longleftrightarrow\quad B_j\cap A_l\ne\varnothing.
\]
Every vertex has at least one incoming and one outgoing edge in this relation.
For an equivalence relation $R$ on $[d]$, define $f(R)$ to be the equivalence
relation generated by all pairs $(j,j')$ for which there are $l,l'$ satisfying
\[
 j\mathrel D l,\qquad j'\mathrel D l',\qquad l\mathrel R l'.
\]
The operator $f$ is monotone for inclusion of equivalence relations.

\begin{lemma}[Finite component stabilization]
\label{lem:through-stabilization}
For $n\geq1$, let $E_n$ relate $j$ and $j'$ when their component occurrences
in the first layer belong to the same connected component of $T^{\circ n}$.
Then
\[
 E_1=\Delta_{[d]},\qquad E_{n+1}=f(E_n),\qquad E_n\subseteq E_{n+1}.
\]
If $E_n=E_{n+1}$ once, all subsequent relations agree.  In particular,
$E_d=E_{d+s}$ for every $s\geq0$.
\end{lemma}

\begin{proof}
The relation for one layer is the identity because the $C_j$ are precisely
the connected components of $T$.  Decompose $T^{\circ(n+1)}$ as its first
layer followed by a power of length $n$.  The latter power connects its
first-layer components according to $E_n$.  Gluing along the seam therefore
connects two components $j,j'$ of the new first layer exactly by chains of
the pairs defining $f(E_n)$.  This proves the recurrence.

The identity relation is contained in every equivalence relation, so
$E_1\subseteq E_2$.  Monotonicity of $f$ now proves the remaining inclusions
inductively.  Equality at consecutive terms is a fixed point of the recurrence
and hence persists.  Every strict inclusion merges at least two equivalence
classes.  There are initially $d$ classes, so there are at most $d-1$ strict
inclusions before stabilization.  Thus $E_d$ is already stable.
\end{proof}

The number of classes of $E_n$ equals the number $\rho(T^{\circ n})$ of
components of $T^{\circ n}$, because every component meets the first layer.
This is where the through-component hypothesis is essential.

\begin{lemma}[The rank criterion for alignment]
\label{lem:through-rank-alignment}
If $S$ is a through tile, then
\[
 \rho(S^{\circ2})=\rho(S)
 \quad\Longleftrightarrow\quad
 S\text{ has aligned component partitions}.
\]
\end{lemma}

\begin{proof}
Let $r=\rho(S)$.  Form the bipartite intersection graph with $r$ vertices
on either side, one for each component of the first or second copy of $S$.
Join a first-copy component to a second-copy component if their seam index
sets intersect.  There are no isolated vertices, and its connected components
are exactly those of $S^{\circ2}$.

Such a bipartite graph has at most $r$ components.  It has exactly $r$ if
and only if each component contains one vertex on each side, so its adjacency
relation is a perfect matching.  Because both families of seam index sets
are partitions, that matching condition says exactly that each output-index
set equals the uniquely matched input-index set.  This is alignment.
\end{proof}

\begin{theorem}[Linear reblocking bound]
\label{thm:through-reblocking}
Every through tile $T$ with $d$ components has an aligned power
$T^{\circ m}$ for some $1\leq m\leq d$; in particular $T^{\circ d}$ is
aligned.  If $m$ is the least such exponent and $r=\rho(T^{\circ m})$, then
\[
 m\leq d-r+1.
\]
For every $n<m$,
$\rho(T^{\circ(n+1)})<\rho(T^{\circ n})$.
The bound $m\leq d$ is attained for every $d\geq2$.
\end{theorem}

\begin{proof}
Lemma~\ref{lem:through-stabilization} gives
$\rho(T^{\circ d})=\rho(T^{\circ2d})$.  Apply
Lemma~\ref{lem:through-rank-alignment} to $S=T^{\circ d}$ to obtain
alignment.

If $\rho(T^{\circ(n+1)})=\rho(T^{\circ n})$, the inclusion
$E_n\subseteq E_{n+1}$ must be equality.  Stabilization then implies
$\rho(T^{\circ2n})=\rho(T^{\circ n})$, and the rank criterion shows that
$T^{\circ n}$ is aligned.  Before the first aligned power, the rank must
therefore drop at every step.  Starting from $d$ and ending at $r$ gives
$m-1\leq d-r$, as required.

For sharpness, use $w=d+1$ interface indices $0,1,\ldots,d$ and the sets
\[
 A_1=\{0,1\},\qquad A_j=\{j\}\ (2\leq j\leq d),
\]
\[
 B_1=\{0,2\},\qquad B_j=\{j+1\}\ (2\leq j<d),\qquad B_d=\{1\}.
\]
Choose any connected graph for each component on its prescribed input and
output ports, with distinct components disjoint.  The seam relation is the
directed cycle $1\to2\to\cdots\to d\to1$ together with the loop $1\to1$.
For $1\leq n\leq d$, the relation $E_n$ has one nonsingleton class
\[
 \{1,d,d-1,\ldots,d-n+2\},
\]
interpreted as $\{1\}$ when $n=1$, and every remaining class is a singleton.
This follows immediately by induction from the recurrence for $f$.
Consequently $\rho(T^{\circ n})=d-n+1$ through exponent $d$, where the
power becomes connected.  An aligned power would have stable rank at all
later powers by the component-chain decomposition.  Thus no exponent below
$d$ is aligned, proving sharpness.
\end{proof}

\subsection{A kernel for all through-component powers}

\begin{corollary}[Polynomial marked-cut kernel after reblocking]
\label{cor:through-kernel}
Let $T$ be a through tile of explicit size
$Q=|V(T)|+|E(T)|$ with $d$ components.  For every fixed $k\geq0$ and
binary exponent $N\geq1$, the open power $T^{\circ N}$ has a
context-preserving marked-signature kernel constructible in polynomial bit
time in the encoded input size and $\log(N+1)$.

Choose the least aligned exponent $m\leq d$, put $S=T^{\circ m}$ and
$r=\rho(S)$, and write $N=qm+t$, $0\leq t<m$.  If $q=0$, retain the
whole power of length $N<m$.  Otherwise use
\[
 K(S,q)\circ T^{\circ t},
\]
where the empty suffix is omitted.  The explicit size is at most
\[
 (a_k+1)r\,mQ+(m-1)Q
 \leq (a_k+2)d^2Q.
\]
At threshold three this gives the concrete bound $6d^2Q$.
\end{corollary}

\begin{proof}
At most $d$ explicit copies of $T$ suffice to find the first aligned power
by serial construction and component analysis.  Its size is at most $mQ$.
The identity $T^{\circ N}=S^{\circ q}\circ T^{\circ t}$ is an exact
rebracketing before any reduction.  Apply the aligned kernel theorem to the
first factor and retain the suffix exactly.  Its contextual form establishes
the marked-signature assertion.  Counting retained records gives the displayed
bound, and binary division by $m\leq d$ plus the polynomial-size construction
gives the bit-time claim.  The case $q=0$ never uses a zero-exponent tile.
\end{proof}

The sharper orbit-sensitive bound is
$ (a_k+1)\sum_{\mathcal O} h_{\mathcal O}Q_{\mathcal O}+(m-1)Q$,
where the cycles and component sizes are those of the explicit tile $S$.
Using the least aligned exponent also gives $rm\leq m(d-m+1)$.
These refinements can substantially reduce the conservative quadratic bound.

\paragraph{Original-source witnesses.}
While explicitly forming $S=T^{\circ m}$, retain for each edge its pair
$(u,e)$ consisting of its zero-based original tile-copy index $0\leq u<m$
and original edge-occurrence identifier $e$.  A lifted cut or mark address
$(v,(u,e))$ in copy $v$ of $S$ then has original copy index
\[
 mv+u.
\]
For a retained suffix edge at suffix index $u$, the original copy index is
$mq+u$.  These formulas compose with the cycle-word run witnesses and
preserve the exact cut cost and marks.  A checker can reconstruct $S$ from
the original tile, replay the finite routing certificate for $S$, and replay
the source coloring using this affine address map.  All copy indices require
only $O(\log(N+1))$ bits.  None of these steps expands the binary exponent.

\paragraph{Cubic terminal use.}
Reblocking preserves the actual original seams, ports, and edges.  If the
original closed flat chain satisfies the cubic local-degree checks, then so
does its exact rebracketing; the subsequent routing kernel preserves those
degrees as before.  The geometric limitation remains unchanged: this is a
graph-query theorem and requires a separately certified ball for its terminal
topological interpretation.

\paragraph{Scope and audit.}
The reblocking argument does not cover a component missing one of the
interfaces.  Such components are not elements of the block-bijection model,
and the rank recurrence above no longer counts every component of the power.
Its extension to those components would require a further proof.
The stabilization claim was exhaustively checked on all $7{,}089$ block
bijection diagrams of widths one through five (respectively $1,3,25,339,6721$
diagrams), using explicit graph composition.  In every tested unaligned
power the next rank strictly dropped.  This audit supports implementation;
the lemmas above supply the proof for arbitrary width.

% Merge these records into the bibliography of the parent article.
